
A language model that achieves high accuracy on MMLU may still hallucinate on a substantial fraction of TruthfulQA prompts, collapse under adversarial perturbations, and emit confidently incorrect answers. These failure modes are not necessarily independent: they may co-vary along a shared latent dimension that capability scores do not capture. In practice, organizations deploying open-weight LLMs rely predominantly on capability benchmarks to select models, an approach that may systematically overlook safety-relevant properties orthogonal to raw knowledge.

The evaluation community has established dedicated benchmarks for calibration (Expected Calibration Error, Brier score), hallucination resistance (TruthfulQA), and adversarial robustness (AdvGLUE, ANLI). What has not been examined at the population level is whether these dimensions share underlying latent structure. If they are genuinely orthogonal, each requires independent measurement; if they co-vary along a latent dimension, a compact metric battery could efficiently identify high-risk models. The gap is analytic rather than measurement-based: the benchmarks exist, but a capability-controlled cross-property correlation matrix computed across a diverse open-weight model population, coupled with factor analysis to test for a coherent latent dimension, has not been reported.

Comprehensive evaluation frameworks such as HELM (Liang et al., 2022), DecodingTrust (Wang et al., 2023), and TrustLLM (Sun et al., 2024) report per-metric scores across many dimensions but do not analyze the factor-analytic structure of cross-metric covariance across model populations. This paper addresses that gap.

The guiding construct is \textit{epistemic reliability}: the degree to which a model's outputs faithfully reflect its internal uncertainty. A model that accurately represents what it does not know may be less prone to hallucination, better calibrated, and more resistant to adversarial perturbations that exploit overconfident decision boundaries. Crucially, this property is posited to be largely orthogonal to raw knowledge as measured by MMLU.

This work introduces the \textbf{YouRA evaluation framework} and reports pipeline validation results under synthetic data conditions. The contributions are:

\begin{enumerate}
\item \textbf{A psychometric methodology for LLM trustworthiness analysis.} Treating models as subjects and benchmark scores as items enables factor analysis and partial correlation approaches that reveal latent structure not visible in per-metric reporting.
\end{enumerate}

\begin{enumerate}
\item \textbf{Pipeline validation results under synthetic conditions.} Under a synthetic score matrix conforming to the hypothesized structure (N=30, 5 metrics), partial Spearman correlations controlling for MMLU are strong (ECE–TruthfulQA\%: ρ=−0.758; ECE–AdvGLUE drop: ρ=−0.719), a single factor explains 72.1\% of shared variance (Tucker's congruence = 1.000), and the survival fraction of 0.943 indicates that controlling for MMLU reduces the calibration–hallucination correlation by 5.7\%.
\end{enumerate}

\begin{enumerate}
\item \textbf{A null result on incremental predictive power.} The composite epistemic predictor achieves LOO-AUC = 0.739 but the incremental advantage over MMLU-only (ΔAUC = 0.051, BCa 95\% CI [−0.194, 0.449]) does not meet the pre-specified threshold and the CI includes zero.
\end{enumerate}

\begin{enumerate}
\item \textbf{A pre-registered analysis pipeline for real-data replication.} The full lm-evaluation-harness pipeline is implemented. Real-data execution (FW1) is identified as the immediate priority before any claims about actual LLM behavior can be made.
\end{enumerate}

All quantitative results reported here are derived from synthetic data. These results demonstrate pipeline correctness under controlled conditions.

